\usection{Abstract}

This thesis asks when two communication-avoiding (CA) mechanisms improve a
Krylov exponential integrator for a three-asset Black--Scholes PDE.
Matrix-powers tiling reduces data movement in sparse operator applications,
while grouped orthogonalization reduces global reductions. Both mechanisms
also add work, so they can lose when redundant arithmetic, halo exchanges, or
local kernel costs outweigh the communication saved.

A dimensionless regime map organizes this comparison. Its vertical coordinate
$\Rv$ compares the Arnoldi working set with last-level-cache capacity, and its
horizontal coordinate $\Rh$ compares reduction latency with useful computation
between reductions. Both are computed before timing from the problem and
independently calibrated machine constants. Their unit lines identify
opportunities, not universal crossovers. Cache geometry, access efficiency,
overlap, and topology still determine whether an implementation recovers the
potential saving.

Neither mechanism pays on the tested CPU. An 11--16$\times$ reduction-count
decrease gives no wall-clock improvement at $\Rh=0.080$, and halo work and
access geometry keep cache-blocked matrix powers below their traffic ceiling.
On a V100, a plane-streamed schedule cuts executed instructions by
2.88--3.45$\times$. A measured runtime rule retains full volume for $n<61$ and
selects plane streaming from $n=61$ when shared memory permits. It improves the
complete solver by 1.248$\times$ geometrically on four selection cases and
1.239$\times$ on four active held-out cases, with worst cases of 1.204$\times$
and 1.182$\times$; four fallback controls remain near parity.

An equal-work distributed control removes an unused recurrence application
from the $s=1$ path. The resulting algorithm-relative ratio is
0.98--1.01$\times$ on one V100 and 1.78--1.80$\times$ on four V100s across two
nodes. The gain is topology-specific: four V100s deliver only 0.89$\times$ the
one-GPU performance at fixed width, and the best four-GPU configuration remains
12--15\% slower than the best one-GPU configuration.

The main contribution is a controlled framework for deciding when
communication avoidance offers an opportunity and explaining why an
implementation does or does not recover it. Independent Johnson and
randomized quasi-Monte Carlo references establish total-price and Greek error,
while measured orthogonality and an independent exponential-action referee
provide the numerical acceptance evidence. The pricing application is used as
a demanding case study; it is not presented as a state-of-the-art pricer.

\textbf{Keywords:} Krylov subspace methods, communication avoidance, matrix exponential, option pricing, performance modeling
